\begin{abstract}
\noindent
Multi-agent LLM coding systems load a static instruction document
(\claudemd or equivalent) into every agent spawn as a cached prefix.
As these documents grow through iterative maintenance, they dominate
first-spawn token cost and human maintenance burden. The na\"{i}ve
response --- lossy compression --- lacks the instruction-following
benchmarks needed to justify it safely.

We propose \emph{frequency-weighted progressive disclosure}: a
conservative trim algorithm that partitions the prefix into sections,
classifies each as \emph{must-resident}, \emph{frequency-deciding}, or
\emph{redundant}, and externalises or drops only content that the
decision rule confirms is cheaper to fetch on demand than to carry
in-prefix on every spawn. No instruction text is rewritten or
compressed; only duplication is removed and infrequently-needed
content is replaced by one-line pointers to permanent external homes.

We apply the algorithm to \blackrim's 708-line, 48\,487-character
\claudemd across three waves of revision. The trim removes 249 lines
(35.2\,\%) and 9\,419 characters (19.4\,%), with the first two waves
delivering a 21.5\,\% line reduction. To verify fidelity, we score a
16-prompt held-out evaluation covering four dimensions (worktree
isolation, delegation discipline, commit path, and merge flow) under
binary pass/fail rubric. The trimmed instructions improve measured
instruction-following by 18.8 percentage points (68.8\,\%
$\to$ 87.5\,\% pass rate at $-21.5$\,\% lines), consistent with the
hypothesis that redundant prose dilutes load-bearing rules. The trim
also reduces amortised input token cost by approximately 17\,\% under
observed cache hit rates. All scripts, CSVs, and evaluation responses
are archived with the paper for full reproducibility.
\end{abstract}
